\begin{textsample}{POS Dim 4 – global\_south\_2024\_07 – Score 21.00 – global\_south\_2024\_07/110144480.txt}  \label{ex:f4_pos_120}
Israeli soldiers gather in the southern Gaza Strip on Wednesday, July 3, 2024.
The Israeli military invited reporters for a tour of Rafah, where the military has been operating since May 6.
DUBAI, United Arab Emirates -- Several officials in the Middle East and the US believe the level of devastation in the Gaza Strip caused by a nine-month Israeli offensive likely has helped push Hamas to soften its demands for a cease-fire \textbf{agreement}.
Hamas over the weekend appeared to drop its longstanding demand that Israel promise to end the war as part of any cease-fire \textbf{deal}.
The sudden shift has raised new hopes for progress in internationally brokered \textbf{negotiations}.
Israeli \textbf{Prime} Minister Benjamin Netanyahu on Sunday boasted that military pressure -- including Israel ' s ongoing two-month offensive in the southern Gaza city of Rafah --"is what has led Hamas to enter \textbf{negotiations}."Hamas, an Islamic militant group that seeks Israel ' s destruction, is highly secretive and little is known about its inner workings.
But in recent internal communications seen by The Associated Press, messages signed @ @ @ @ @ @ @ @ @ @ ' s \textbf{exiled} political leadership to \textbf{accept} the cease-fire \textbf{proposal} pitched by US \textbf{President} Joe Biden.
The messages, shared by a Middle East official familiar with the ongoing \textbf{negotiations}, described the heavy losses Hamas has suffered on the battlefield and the dire conditions in the war-ravaged territory.
The official spoke on condition of anonymity to share the contents of internal Hamas communications.
It was not known if this internal pressure was a factor in Hamas ' flexibility.
But the messages indicate divisions within the group and a readiness among top militants to reach a \textbf{deal} quickly, even if Hamas ' top official in Gaza, Yahya Sinwar, may not be in a rush.
Sinwar has been in hiding since the war erupted last October and is believed to be holed up in a tunnel deep underground.
US officials declined to comment on the communications.
But a person familiar with Western intelligence who spoke on the condition of anonymity to discuss the sensitive matter said the group ' s leadership understands its forces have suffered heavy losses and @ @ @ @ @ @ @ @ @ @.
Two US officials say the Americans are aware of internal divisions within Hamas and that those divisions, the destruction in Gaza or pressure from \textbf{mediators} Egypt and Qatar could have been factors in the militant group softening its demands for a \textbf{deal}.
The US officials spoke on the condition of anonymity to discuss the Biden administration ' s view of the current situation.
The Middle Eastern official shared details from two internal Hamas communications, both written by senior officials inside Gaza to the group ' s \textbf{exiled} leadership in Qatar, where Hamas ' supreme leader, Ismail Haniyeh, is based.
The communication suggested that the war had taken a toll on Hamas fighters, with the senior figures urging the militant ' s political wing abroad to \textbf{accept} the \textbf{deal} despite Sinwar ' s reluctance.
Hamas spokesperson Jihad Taha dismissed any \textbf{suggestions} of divisions within the group."The movement ' s position is unified and is crystallized through the organizational framework of the leadership,"he said.
The intelligence official showed The Associated Press a transcript @ @ @ @ @ @ @ @ @ @ specific details about how the information was obtained, or the raw form of the communications.
The official said the communications took place in May and June and came from multiple senior officials inside the group ' s military wing in Gaza.
The messages acknowledged Hamas fighters had been killed and the level of devastation to the Gaza Strip wrought by the Israeli campaign in the enclave.
They also suggest that Sinwar either isn't fully aware of the toll of the fighting or isn't fully communicating it to those negotiating outside of the territory.
It was not known whether Haniyeh or any other top officials in Qatar had responded.
Israeli officials declined to comment on the communications.
Egypt and Qatar also had no immediate comment.
Egypt and Qatar have been working with the United States to broker a cease-fire and end the devastating nine-month war.
After months of fits and starts, \textbf{talks} \textbf{resumed} last week and are scheduled to continue in the coming days.
A \textbf{deal} is still not \textbf{guaranteed}.
Netanyahu ' s @ @ @ @ @ @ @ @ @ @."The US officials said they are cautiously \textbf{optimistic} about the prospects for a cease-fire based on the latest developments, but stressed that numerous efforts had looked promising only to fall through.
Still, the sides appear closer to a \textbf{deal} than they have been in months.
Israel launched the war in Gaza after Hamas ' October attack in which militants stormed into southern Israel, killed some 1,200 people -- mostly civilians -- and abducted about 250.
Israel says Hamas is still holding about 120 hostages -- about a third of them thought to be dead.
Since then, the Israeli air and ground offensive has killed more than 38,000 people in Gaza, according to the territory ' s Health Ministry, which does not distinguish between combatants and civilians.
The offensive has caused widespread devastation and a humanitarian crisis that has left hundreds of thousands of people on the brink of famine, according to international officials.
Both Hamas and \textbf{Egyptian} officials confirmed Saturday that Hamas has dropped a key demand that Israel commit up front @ @ @ @ @ @ @ @ @ @ demand, leaving the \textbf{talks} stalled for months.
Instead, the officials, speaking on condition of anonymity to discuss ongoing \textbf{negotiations}, said the \textbf{phased} \textbf{deal} would start with a \textbf{six-week} cease-fire during which older, sick and female hostages would be \textbf{released} by Hamas in \textbf{exchange} for hundreds of Palestinian \textbf{prisoners}. \textbf{Talks} on a broader \textbf{deal}, including an end to the war, would only begin during this \textbf{phase}, they said.
Netanyahu has vowed to keep fighting until Israel destroys Hamas ' military and governing capabilities, even if hostages are \textbf{freed}.
The Associated Press writers Aamer Madhani and Matthew Lee in Washington contributed to this report.

% matched lemmas: accept, agreement, deal, egyptian, exchange, exile, free, guarantee, mediator, negotiation, optimistic, phase, president, prime, prisoner, proposal, release, resume, six-week, suggestion, talk
\end{textsample}
